% ==============================================================================
% secoes/02_fundamentacao.tex - SEÇÃO II: FUNDAMENTAÇÃO TEÓRICA
% Tema: Federico Faggin e o Primeiro Microprocessador Monolítico (Intel 4004)
% ==============================================================================

\section{Fundamentação Teórica}
\label{sec:fundamentacao}

% [DIRETIVA DE ESCRITA: INTRODUÇÃO DA FUNDAMENTAÇÃO TEÓRICA]
% - Estabelecer a transição dos transistores bipolares para a tecnologia MOS e
%   o impasse enfrentado pelos processos de fabricação até o final dos anos 1960.

\subsection{A Tecnologia MOS de Porta Metálica vs. Porta Autoalinhada (SGT)}
\label{subsec:mos_sgt}

% [DIRETIVA DE ESCRITA: LIMITAÇÕES DA PORTA DE ALUMÍNIO E A DESCOBERTA DO SGT]
% - Descrever a física dos transistores MOS com porta de alumínio (metal-gate MOS):
%   a necessidade de tolerâncias de tolerância de alinhamento mecânico das máscaras fotolitográficas,
%   resultando em expressiva sobreposição geométrica (overlap) entre a porta e as regiões de fonte e dreno.
% - Explicar como Federico Faggin, em 1968 no centro de P&D da Fairchild Semiconductor,
%   desenvolveu a Silicon Gate Technology (SGT) \cite{faggin1970}:
%   a utilização de polissilício termicamente estável depositado antes da dopagem das regiões de
%   fonte e dreno, permitindo que a própria porta atue como máscara autoalinhada.
% - Detalhar a fabricação do primeiro chip comercial do mundo em SGT: o Fairchild 3708
%   (multiplexador analógico de 8 bits com lógica decodificadora), comprovando a viabilidade industrial do processo.

\subsection{Vantagens Físicas e Métricas da Silicon Gate Technology}
\label{subsec:vantagens_sgt}

% [DIRETIVA DE ESCRITA: PROPRIEDADES FÍSICAS E GANHOS DE DESEMPENHO DO SGT]
% - Analisar detalhadamente as três vantagens fundamentais introduzidas pelo SGT:
%   1. Minimização da Capacitância Parasita de Sobreposição (Cov): ao eliminar o overlap mecânico,
%      a capacitância Miller entre porta e dreno é reduzida drasticamente, aumentando a velocidade
%      de comutação intrínseca dos transistores em até 5 vezes frente ao metal-gate equivalente.
%   2. Redução e Estabilidade da Tensão de Limiar (Vth): a diferença de função de trabalho
%      (work function) entre o silício dopado e o substrato resulta em tensões de limiar mais baixas
%      e confiáveis, viabilizando a compatibilidade de alimentação com níveis lógicos TTL (Transistor-Transistor Logic).
%   3. Camada Condutora Adicional para Roteamento: o polissilício condutor atua como um segundo
%      nível de interconexão abaixo do alumínio, permitindo cruzamentos ortogonais sem curtos-circuitos
%      e aumentando drasticamente a densidade de integração em layouts de lógica aleatória (random logic).

\subsection{O Ecossistema do Sistema MCS-4 (Micro Computer Set)}
\label{subsec:sistema_mcs4}

% [DIRETIVA DE ESCRITA: O CONJUNTO MCS-4 E A CPU 4004]
% - Apresentar a arquitetura distribuída do conjunto de chips MCS-4 projetado por Faggin:
%   * Intel 4001: Memória ROM de 2048 bits com portas de I/O paralelas integradas;
%   * Intel 4002: Memória RAM de 320 bits com portas de saída e registradores dinâmicos;
%   * Intel 4003: Registrador de deslocamento de 10 bits para expansão serial de I/O;
%   * Intel 4004: A Unidade Central de Processamento (CPU) de 4 bits contendo cerca de 2.300 transistores
%     PMOS integrados em uma pastilha de apenas 12 mm², encapsulada em padrão DIP de 16 pinos.
% - Enfatizar a genialidade do encapsulamento compacto: o uso de um barramento multiplexado de endereço
%   e dado de 4 bits permitiu acomodar um computador completo de uso geral em pinagem restrita,
%   viabilizando custos de embalagem extremamente reduzidos.
